\section{Introduction}\label{sec:introduction}

Let \(\Omega\subset\R^2\) be a bounded strictly convex open set whose
boundary is a \(C^\infty\) simple closed regular curve with positive
curvature. Orient the boundary counterclockwise, write its arclength
coordinate as \(s\in\R/L\mathbb Z\), and let \(\delta\in(0,\pi)\) be the
angle of the outgoing velocity measured counterclockwise from the
positive tangent. Reflection at a boundary point with outward unit
normal \(n\) sends an incoming velocity \(v\) to
\(v-2(v\cdot n)n\). The resulting billiard map \(T\) acts on
\[
M=(\R/L\mathbb Z)\times(0,\pi),
\]
and \(\delta\downarrow0\) is the grazing end considered below.
Write \(\T=\R/2\pi\mathbb Z\) for the angular circle.

\begin{theorem}[Phase-foliation rigidity]\label{thm:main}
Suppose a \(T\)-invariant open annulus \(U\subset M\) contains
\((\R/L\mathbb Z)\times(0,\eta)\) for some \(\eta>0\).
Suppose there is a homeomorphism
\[
H:\mathbb S^1\times(0,1)\longrightarrow U
\]
whose leaves are essential simple closed curves, each individually
invariant under \(T\). Then \(\Omega\) is an ellipse: there are
\(c\in\R^2\) and a real symmetric positive-definite matrix \(Q\) such that
\[
\Omega=\{x\in\R^2:(x-c)^TQ(x-c)<1\}.
\]
Circles are included.
\end{theorem}

For an elliptical table, the confocal elliptic caustics give such a
grazing foliation \citep{Lazutkin1973}; for a circular table, the
corresponding caustics are concentric circles.
Only continuity of the leaves and of their transverse dependence is
assumed. The foliation need not extend to the rest of the phase cylinder,
and the width \(\eta\) may be arbitrarily small. The main issue is that a
continuous essential phase curve need not initially be given as the
tangent lines of a smooth convex caustic.

The complete companion article \emph{Rigidity of Smooth Billiards with a
Continuous Caustic Collar} \citep{NB} supplies three analytic inputs and
the final physical-collar rigidity theorem. Its stationary-mesh and
Fourier-row estimates apply to arbitrary normalized actions; its buffered
division lemma is an analytic estimate. Their precise hypotheses are
recalled at the points of use. Here we establish the additional dynamical
and geometric steps needed to apply them to a merely continuous phase
foliation. In particular, we construct the physical collar before invoking
the companion's ellipse theorem.

\subsection{Context and significance}

The Birkhoff conjecture asks whether an integrable strictly convex planar
billiard must be an ellipse. Birkhoff developed the convex billiard as a
dynamical model \citep{Birkhoff1927}; Poritsky's early printed formulation
also studied compatible string constructions, under hypotheses stronger
than a foliation alone \citep{Poritsky1950}. The distinction between a
full foliation and a family with gaps is essential. Lazutkin's theorem
guarantees a Cantor family of near-boundary invariant curves for
sufficiently smooth positively curved tables, rather than a full foliation
\citep{Lazutkin1973}. Marvizi and Melrose constructed a grazing first
integral modulo a flat error, not an exact invariant foliation
\citep{MarviziMelrose1982}. Positive curvature is substantive: Mather
excluded grazing caustics when boundary curvature vanishes
\citep{Mather1982}.

Stronger or different integrability hypotheses give complementary
rigidity results. Bialy proved that foliation of the entire phase
cylinder by noncontractible invariant circles forces a circle
\citep{Bialy1993}. Avila, De Simoi, and Kaloshin proved local rigidity
near small-eccentricity ellipses \citep{ADK2016}, and Kaloshin and
Sorrentino treated perturbations of ellipses of arbitrary eccentricity
\citep{KS2018}. Bialy and Mironov allowed a continuous phase foliation
under central symmetry when the foliated region reaches the invariant
four-periodic curve \citep{BialyMironov2022}. Koval's local strong
result treats caustics at all rational rotations in an arbitrarily small
interval above zero, under closeness to almost every ellipse and
high-regularity bounds
\citep{Koval2026}. These results impose geometric restrictions or require
invariant curves extending beyond an arbitrarily thin grazing annulus.

Birkhoff's invariant-curve theorem for twist maps explains why graphs are
the natural intermediate objects \citep{Arnaud2009}. Here the initial
leaves are only continuous essential curves, and the entire annulus is
foliated. We therefore prove graphness directly from both billiard twist
signs and finite-area recurrence. The full collar then supplies a wholly
periodic graph at every sufficiently small rational rotation. The
stationary actions of those graphs yield high-frequency equations for the
boundary. First-order Fourier obstructions for rational caustics in
perturbations of the circle were developed by \citet{RamirezRos2006};
quantitative constraints near ellipses appear in \citet{ADK2016}.
The analytic inversion here uses the companion's arbitrary-input estimates
without a perturbative assumption.

Together with the physical-collar rigidity theorem of \cite{NB},
Theorem~\ref{thm:main} resolves the near-boundary Birkhoff conjecture
positively in the continuous phase-foliation formulation above.
The additional point established here is the passage from merely
continuous invariant phase leaves to the smooth physical caustic
hypothesis. This passage requires neither a symmetry assumption nor
closeness to a circle or ellipse.

\subsection{The regularity reduction}

A \emph{smooth convex caustic} is a \(C^\infty\) simple closed regular
curve in \(\Omega\), bounding a strictly convex set, such that a billiard
segment tangent to it remains tangent after reflection at either
endpoint. A \emph{continuous physical caustic collar} is a relatively
open set \(A\subset\overline\Omega\), containing \(\partial\Omega\),
together with a homeomorphism
\[
H_0:\mathbb S^1\times[0,1)\longrightarrow A
\]
whose zero leaf is \(\partial\Omega\) and whose positive leaves are smooth
convex caustics. This is the hypothesis of the following companion theorem.

\begin{theorem}[Physical-collar rigidity,
  {\citep[Definition~1.1 and Theorem~1.2]{NB}}]
\label{thm:nb-rigidity}
A bounded strictly convex planar domain with \(C^\infty\) regular boundary
of positive curvature and a continuous physical caustic collar is an
ellipse.
\end{theorem}

The core reduction proved below is stronger than the smoothness needed
for this application.

\begin{theorem}[Regularity of a phase-foliated collar]
\label{thm:collar-regularity}
Under the hypotheses of Theorem~\ref{thm:main}, the boundary is real
analytic. There are \(\lambda_0>0\) and a jointly real-analytic map
\[
\Gamma:\mathbb T\times[0,\lambda_0]\longrightarrow\R^2
\]
such that \(\Gamma(\cdot,0)\) parametrizes \(\partial\Omega\), every
curve \(\Gamma(\cdot,\lambda)\) for \(0<\lambda\le\lambda_0\) is a regular
strictly convex caustic, and the restriction
to \(\mathbb T\times[0,\lambda_0)\) is a homeomorphism onto a relatively
open collar of \(\partial\Omega\) in \(\overline\Omega\).
Here real analyticity at the parameter endpoints means extension to a
neighborhood.
\end{theorem}

Theorem~\ref{thm:collar-regularity} uses the mesh and Fourier estimates
from \cite{NB}. Its conclusion is proved directly from the phase
foliation, without assuming the physical collar that appears in
Theorem~\ref{thm:nb-rigidity}.

\subsection{Structure of the argument}

The first step is topological and dynamical. In oriented-line
coordinates the billiard has two strict twist inequalities. They impose
one-way barriers on the lifted argument of successive angular
separations. On an interval of equal rotation numbers, area recurrence
can be applied on a suitable finite-measure part of the simultaneous
cover of the product annulus. The resulting contradiction excludes
vertical pairs. A separate rational-passage argument gives graph leaves
densely without first assuming monotonicity of the rotation number.
Limits then give graphs everywhere. Area recurrence also rules out
rational plateaus, making each sufficiently small rational graph unique
and wholly periodic. This graph argument is useful independently of
the analytic estimates that follow.

The second step obtains boundary analyticity from stationary polygonal
chains. Because the initial line lies on an invariant graph, a chain
whose angle returns must also return its momentum. Its stationary action
is therefore constant in the initial point. The associated
high-frequency equations have derivative equal to the high-mode
projection plus a small operator. A Newton construction on decreasing
complex strips produces an analytic solution, and a real contraction
identifies it with the original smooth boundary data.

The third step is a joint analytic construction in the angular and
rotation variables. At rational resonances it uses both residual values
and first parameter jets. The first jets control the coefficients left
undetermined by one Fourier division, allowing a second division and a
quadratic Newton step. The polynomial loss exponent is fixed before the
order of formal initialization. These mechanisms come from the analytic
machinery of \cite{NB}; we give the argument with its weaker geometric
input explicit.

Finally, the invariant line graphs define analytic support functions
with positive curvature radius. Their support decreases strictly to
first order in the squared translation parameter, so their envelopes
are nested convex caustics filling a physical collar.
Figure~\ref{fig:phase-to-caustic} illustrates this final passage from
an analytic line graph to its support envelope.
This proves Theorem~\ref{thm:collar-regularity} and makes
Theorem~\ref{thm:nb-rigidity} applicable.

Section~\ref{sec:geometry} proves the graph and rational-periodicity
statements. Section~\ref{sec:analytic} establishes analytic boundary
regularity. Section~\ref{sec:joint} constructs the joint parametrization,
and Section~\ref{sec:caustics} proves the physical collar and concludes
the rigidity argument.

\paragraph{Conventions.}
All angular coordinates are real lifts of \(2\pi\)-periodic variables.
Translation numbers are measured in radians; the corresponding rotation
number is the translation number divided by \(2\pi\).
Constants may increase between estimates, with their dependence and
order of choice specified when necessary.
A periodic holomorphic function is \emph{real symmetric} if
\(f(\bar z)=\overline{f(z)}\); the same convention applies in two variables.
